%==============================================================================
%  MATH 231 -- Linear Algebra I (Queens College, Fall 2026)
%  PLAN / NOT FINALIZED -- How the rest of the semester's material is split
%  into lecture documents: the matrix unit, Parts 2-11 (Midterm 2); then, for
%  the final, the eigen unit and the linear-maps unit (linear maps,
%  composition, change of basis, the SVD), which closes the course.
%
%  SOURCE.  The topic lists and their ORDER are the instructor's (2026-10-25,
%  given in three messages: the matrix unit, then the eigen unit, then the
%  closing topics), including the derivations the instructor specified step
%  by step (the orthogonal-matrix lemma, the area of a parallelogram, the
%  rotation matrix from the unit circle, least squares from the scalar
%  problem, rank(A'A) = rank(A) from rank-nullity).  The order is kept
%  exactly.  What this document
%  adds is (1) where each topic lands -- which Part, which section number --
%  (2) the cross-references into material students already have, and (3)
%  anything proposed beyond the list, tagged [proposed] so it can be struck.
%  Points where the list and the mathematics disagree are tagged [D1]-[D12]
%  and collected in the "Decisions" section.
%
%  WHEN WRITING A PART: read its entry here, the matching Decisions, and the
%  Notation section first.  Each Part's header comment should carry over the
%  decisions it depends on, as Part 1's header does.
%
%  AUDIENCE: the instructor.  The lecture documents themselves stay free of
%  instructor-facing apparatus (pacing, flags); this file is where those live.
%
%  Compile with tectonic or pdflatex (standard TeX Live packages only).
%==============================================================================
\documentclass[11pt]{article}

\usepackage[margin=1in]{geometry}
\usepackage{amsmath,amssymb}
\usepackage[dvipsnames]{xcolor}
\usepackage{enumitem}
\usepackage{booktabs}
\usepackage{array}
\usepackage{longtable}
\usepackage[hidelinks]{hyperref}
\usepackage{titlesec}
\usepackage{needspace}

\titleformat{\section}{\Large\bfseries\sffamily\color{RoyalBlue}}{\thesection.}{0.5em}{}
\titlespacing*{\section}{0pt}{18pt}{8pt}
\titleformat{\subsection}{\large\bfseries\sffamily}{}{0em}{}
\titlespacing*{\subsection}{0pt}{14pt}{4pt}

\setlength{\emergencystretch}{3em}
\setlist{itemsep=2pt,topsep=3pt}

%--------------------------------------------------------------- notation ----
\newcommand{\R}{\mathbb{R}}
\newcommand{\vc}[1]{\mathbf{#1}}
\newcommand{\mt}[1]{\mathbf{#1}}
\newcommand{\T}{^{\mathsf{T}}}
\newcommand{\norm}[1]{\lVert #1\rVert}
\newcommand{\colsp}[1]{\operatorname{col}[#1]}
\newcommand{\rowsp}[1]{\operatorname{row}[#1]}
\newcommand{\nullsp}[1]{\operatorname{null}[#1]}
\DeclareMathOperator{\rank}{rank}
\DeclareMathOperator{\nullity}{nullity}
\DeclareMathOperator{\tr}{tr}
\DeclareMathOperator{\adj}{adj}
\DeclareMathOperator{\spanop}{span}
\DeclareMathOperator{\diagop}{diag}
\newcommand{\proj}[2]{\operatorname{proj}_{#1}(#2)}

%------------------------------------------------------------------ tags -----
\newcommand{\prop}{\,{\color{ForestGreen}\sffamily\scriptsize[proposed]}}
\newcommand{\see}[1]{\,{\color{BrickRed}\sffamily\scriptsize[#1]}}

% \partplan{number}{title}{file stem}{sections}{meetings}
\newcommand{\partplan}[5]{%
  \needspace{6\baselineskip}%
  \subsection{Part #1 \textbullet\ #2}
  \noindent{\small\sffamily File \texttt{#3} \textbullet\ #4 \textbullet\
  about #5 meetings}\par\smallskip}
% \sect{number}{title}
\newcommand{\sect}[2]{\par\smallskip\noindent\textbf{\S#1\ \ #2.}\ \ignorespaces}
\newcommand{\fits}{\par\smallskip\noindent\textit{Homework fit.}\ \ignorespaces}

\title{\vspace{-1.2cm}\textbf{\sffamily MATH 231 -- Linear Algebra I}\\[4pt]
       {\large\sffamily Plan for the Rest of the Term: the Matrix, Eigen,
        and Linear-Maps Units}\\[2pt]
       {\normalsize\sffamily Fall 2026 \textbullet\ working draft --- not finalized}}
\author{}
\date{\normalsize\sffamily Drafted October 25, 2026}

\begin{document}
\maketitle
\vspace{-1.0cm}

\begin{center}
\fcolorbox{BrickRed!50}{BrickRed!7}{%
\begin{minipage}{0.93\textwidth}
\small\sffamily
\textbf{This document is not finalized.} It records how the material is
planned to be split into lecture documents. Nothing in it has been lectured or
announced. Where it and an official course document disagree, the official
document is correct.
\end{minipage}}
\end{center}

%==============================================================================
\section{The split at a glance}
%==============================================================================
The topics, in the order given, become ten new Parts of the matrix unit,
Midterm~2's material, ending with Gram--Schmidt and $QR$. Two shorter units
follow for the final: the eigen unit, then the linear-maps unit, which closes
the course with composition, change of basis, and the SVD. Matrix-unit section
numbers continue across Parts as they do in the vector and computational units,
so Part~2 opens at \S10; each new unit restarts at \S1. Each Part ends with a
Looking-ahead into the next, and each one collects or plants a specific debt,
noted under its entry.

\medskip
{\renewcommand{\arraystretch}{1.3}
\small
\begin{longtable}{@{}l p{4.3cm} l p{6.4cm} c@{}}
\toprule
\textbf{Part} & \textbf{Title} & \textbf{\S} & \textbf{Topics} & \textbf{Mtgs} \\
\midrule
\endfirsthead
\toprule
\textbf{Part} & \textbf{Title} & \textbf{\S} & \textbf{Topics} & \textbf{Mtgs} \\
\midrule
\endhead
2  & Column Space, Rank, and the Determinant & 10--14 &
     $\colsp{A}$, $\rowsp{A}$, rank; a matrix as a map of the plane; the
     $2\times2$ determinant, its effect on the plane, and rank & 1.5 \\
3  & Matrices That Move the Plane & 15--21 &
     scaling (uniform, non-uniform), translation, shear, first rotations,
     reflections, projections; length-preserving matrices; orthogonal
     matrices and what they do & 2 \\
4  & Area, and the Rotation Matrix & 22--25 &
     properties of the $2\times2$ determinant; $|\det\mt{A}|$ as an area;
     $\mt{R}_\theta$ from the unit circle; rotations in $\R^{3}$ & 1--1.5 \\
5  & Projection Matrices and the Direct Sum & 26--29 &
     $\mt{P}_{\vc{u}}$ from the projection formula; idempotence; the
     complementary projection; the direct sum & 1--1.5 \\
6  & The Null Space and Rank--Nullity & 30--34 &
     full rank and rank deficiency, maps between spaces, invariant subspaces;
     null space, kernel, nullity; rank--nullity;
     $\rank(\mt{A}\T\mt{A})=\rank(\mt{A})$; the $2\times2$ inverse & 2 \\
7  & The Determinant in General & 35--39 &
     minors, cofactors, Laplace expansion, principal minors; properties;
     triangular, diagonal, orthogonal & 1.5 \\
8  & Solving Linear Systems & 40--45 &
     row operations, back/forward substitution, Gaussian and Gauss--Jordan
     elimination, (reduced) row-echelon form, homogeneous, over- and
     underdetermined systems & 2 \\
9  & Least Squares & 46--49 &
     the scalar problem; the column space; normal equations; the
     vector-calculus derivation; fitting a line & 1.5 \\
10 & Cramer's Rule, the Inverse, and the Pseudo-inverse & 50--54 &
     Cramer's rule; the $n\times n$ inverse; the adjugate; properties;
     Moore--Penrose & 2 \\
11 & Gram--Schmidt and the $QR$ Decomposition & 55--58 &
     Gram--Schmidt in matrix form; $\mt{A}=\mt{Q}\mt{R}$;
     $\colsp{Q}=\colsp{A}$; $QR$ for least squares & 1.5 \\
\midrule
E1 & Eigenvalues and the Characteristic Polynomial & 1--4 &
     eigenpairs; $\det(\mt{A}-\lambda\mt{I})$; the $2\times2$ formula; trace
     and determinant; Abel--Ruffini & 1.5 \\
E2 & Diagonalization and the Spectral Theorem & 5--7 &
     $\mt{V}\Lambda\mt{V}^{-1}$; symmetric matrices; spectral radius & 1.5 \\
E3 & Structure and Location & 8--9 &
     structured matrices and their eigenvalues; Gershgorin & 1 \\
\midrule
L1 & Linear Maps and Composition & 1--3 &
     linear maps; the matrix of a map; matrix multiplication derived from
     composition & 1 \\
L2 & Change of Basis & 4--5 &
     coordinates in a basis; $\mt{V}^{-1}\mt{A}\mt{V}$; similarity &
     1 \\
L3 & The Singular Value Decomposition & 6--8 &
     the circle to an ellipse; $\mt{U}\Sigma\mt{V}\T$ from
     $\mt{A}\T\mt{A}$; what it gives & 1.5 \\
\bottomrule
\end{longtable}}

\noindent File stems:
\begin{itemize}[leftmargin=1.4em]
  \item \texttt{MATH231\_matrices\_<n>\_\ldots} in \texttt{lectures/matrix-unit/};
  \item \texttt{MATH231\_eigen\_<n>\_\ldots} in a new \texttt{lectures/eigen-unit/};
  \item \texttt{MATH231\_linmaps\_<n>\_\ldots} in a new
        \texttt{lectures/linear-maps-unit/}.
\end{itemize}
Meeting counts are estimates at the pace so far. \S\ref{sec:pacing} sets them
against the calendar, and they do not fit as they stand.

%==============================================================================
\section{The matrix unit, Part by Part}
%==============================================================================

%------------------------------------------------------------------------------
\partplan{2}{Column Space, Row Space, Rank, and the $2\times2$ Determinant}
  {matrices\_2\_column\_space\_rank\_and\_the\_determinant}{\S10--\S14}{1.5}
%------------------------------------------------------------------------------
\noindent\textit{Opens by} picking up Part~1's last question. $\mt{V}\vc{x}=\vc{u}$
is stated but unsolved, and \S2.6 P5 produced a system with no solution. Before
asking \emph{how} to solve (Part~8), ask \emph{when} a solution exists. The
column space answers that. \see{D1}

\sect{10}{The column space} $\colsp{A}=\{\mt{A}\vc{x}:\vc{x}\in\R^{n}\}$, which
by the column reading (\S2.5) is the span of the columns; a subspace of
$\R^{m}$ (vector unit \S19.2--\S19.3). $\mt{A}\vc{x}=\vc{b}$ is solvable exactly
when $\vc{b}\in\colsp{A}$. Revisit \S2.6 P5: $[\,3,1\,]$ is off the line
spanned by $[\,1,2\,]$. The ``no solution'' outcome of the Oct~1 review,
explained.

\sect{11}{The row space} $\rowsp{A}=$ the span of the rows $=\colsp{A\T}$, a
subspace of $\R^{n}$. For $m\neq n$ the two live in different spaces; a
$2\times3$ example with $\colsp{A}\subseteq\R^{2}$ and
$\rowsp{A}\subseteq\R^{3}$. The row reading (\S2.3) is the row space at work.

\sect{12}{Rank} $\rank(\mt{A})=\dim\colsp{A}$, the largest number of
independent columns; $\rank(\mt{A})\le\min(m,n)$. Examples by inspection: the
zero matrix, $\mt{I}_n$, $\begin{bmatrix}1&2\\2&4\end{bmatrix}$, a $3\times3$
whose third column is the sum of the first two, and the outer product
$\vc{u}\vc{v}\T$ (rank one, column space $\spanop\{\vc{u}\}$)\prop. Row rank equals column
rank: checked here for $2\times2$ only, proved in general in \S33.\see{D4}

\sect{13}{A matrix moves the plane} $\vc{x}\mapsto\mt{A}\vc{x}$ as a function
$\R^{n}\to\R^{m}$. Linearity,
$\mt{A}(\alpha\vc{x}+\beta\vc{y})=\alpha\mt{A}\vc{x}+\beta\mt{A}\vc{y}$, from
M2 and M4 (\S8.4). Conversely, a linear map is fixed by where it sends
$\vc{e}_1,\dots,\vc{e}_n$. \emph{Column $j$ of $\mt{A}$ is where $\vc{e}_j$
goes.} This is the first-principles tool Part~3 runs on. The unit square goes
to the parallelogram on the columns, and the grid to a parallelogram grid. The
origin stays put; lines go to lines. Keep this to matrices acting on
$\R^{n}$. Linear maps in general, and the derivation of the product rule from
composition, belong to the linear-maps unit (L1).\see{D11}

\sect{14}{The $2\times2$ determinant} $\det\mt{A}=ad-bc$ is the number
$D=u_1v_2-u_2v_1$ of the vector unit's \S18.5, whose Looking-ahead promised it a
name. That debt is collected here. What it says about the map:
$|\det\mt{A}|$ is the factor by which areas scale (stated here, derived in
\S23), and its sign says whether orientation is kept or flipped.
$\det\mt{A}=0$ holds exactly when the columns are dependent, i.e.\
$\rank\mt{A}<2$, i.e.\ the plane collapses. At rank~1 it collapses onto the line
$\colsp{A}$, and a whole line of vectors lands on $\vc{0}$ (planted for \S31). At
rank~0 everything goes to the origin. One table: rank, determinant, picture.

\fits Catalog Problem~4 (Homework~3 F.1, ``why moving the corners is enough'')
uses exactly \S13.

%------------------------------------------------------------------------------
\partplan{3}{Matrices That Move the Plane}
  {matrices\_3\_matrices\_that\_move\_the\_plane}{\S15--\S21}{2}
%------------------------------------------------------------------------------
\sect{15}{Building a matrix from what it does} The recipe: decide geometrically
where $\vc{e}_1$ and $\vc{e}_2$ go; those are the columns (\S13). Every matrix
below is built that way, then read back from its structure (diagonal,
triangular, symmetric, \dots).

\sect{16}{Scaling} Uniform, $\alpha\mt{I}$: every vector scaled by $\alpha$, every
direction kept, $\det=\alpha^{2}$. Non-uniform, $\diagop(s_1,s_2)$: the axes
scaled separately, so a circle becomes an ellipse, $\det=s_1s_2$, and a negative
entry also flips. Structure: diagonal (\S4.2) means the axes are kept.

\sect{17}{Translation} \see{D2} $\vc{x}\mapsto\vc{x}+\vc{t}$ is not linear, and
no $2\times2$ matrix performs it: $\mt{A}\vc{0}=\vc{0}$ for every $\mt{A}$. It is
done instead in homogeneous coordinates, $[\,x,y\,]\mapsto[\,x,y,1\,]$, by
$\begin{bmatrix}1&0&t_1\\0&1&t_2\\0&0&1\end{bmatrix}$. Combining translation
with the other motions as $3\times3$ products\prop.

\sect{18}{Shear} $\begin{bmatrix}1&k\\0&1\end{bmatrix}$: $\vc{e}_1$ fixed,
$\vc{e}_2\mapsto[\,k,1\,]$, each horizontal line slides in proportion to its
height. The vertical shear is the transpose. Structure: triangular with $1$s on
the diagonal; $\det=1$, so the shape changes and the area does not.

\sect{19}{Rotation, reflection, projection: first examples} Rotation by
$90^\circ$, $\mt{J}=\begin{bmatrix}0&-1\\1&0\end{bmatrix}$, and by $180^\circ$,
$-\mt{I}$; reflection across the $x$-axis, $\diagop(1,-1)$, and across $y=x$,
$\begin{bmatrix}0&1\\1&0\end{bmatrix}$; projection onto the $x$-axis,
$\diagop(1,0)$, with rank~1 and $\det=0$. Each is built from where $\vc{e}_1,\vc{e}_2$ go. A
general angle waits for \S24, a general line for \S26.

\sect{20}{Matrices that preserve length} The question: which
$\mt{A}\in\R^{n\times n}$ satisfy $\norm{\mt{A}\vc{x}}_2=\norm{\vc{x}}_2$ for all
$\vc{x}$? (Start with $n=2$.)
\begin{itemize}[leftmargin=1.4em]
  \item \textbf{Lemma.} If $\mt{M}\in\R^{n\times n}$, $\mt{M}=\mt{M}\T$, and
        $\vc{x}\T\mt{M}\vc{x}=0$ for all $\vc{x}\in\R^{n}$, then $\mt{M}=\mt{0}$.
        \emph{Proof:} $\vc{x}=\vc{e}_i$ gives $m_{ii}=0$; then
        $\vc{x}=\vc{e}_i+\vc{e}_j$ gives
        $m_{ii}+m_{ij}+m_{ji}+m_{jj}=2m_{ij}=0$, using symmetry.
  \item \textbf{Remark}\prop: symmetry cannot be dropped. $\mt{J}$ of \S19 has
        $\vc{x}\T\mt{J}\vc{x}=0$ for every $\vc{x}$, and $\mt{J}\neq\mt{0}$.
        (It is skew-symmetric, Homework~3 B.3.)
  \item \textbf{Theorem.} $\norm{\mt{A}\vc{x}}^{2}=\vc{x}\T\mt{A}\T\mt{A}\vc{x}$
        (\S7.3 with M6) and $\norm{\vc{x}}^{2}=\vc{x}\T\vc{x}$, so
        $\vc{x}\T(\mt{A}\T\mt{A}-\mt{I})\vc{x}=0$ for all $\vc{x}$. The matrix
        $\mt{A}\T\mt{A}-\mt{I}$ is symmetric (M6, \S6.4), so the lemma gives
        $\mt{A}\T\mt{A}=\mt{I}$.
  \item \textbf{Orthonormal columns.} The $(i,j)$ entry of $\mt{A}\T\mt{A}$ is
        $\vc{a}_i\T\vc{a}_j$, so $\vc{a}_i\T\vc{a}_j=\delta_{ij}$, column by column.
        \see{D9: restate $\delta_{ij}$}
  \item \textbf{Definition.} $\mt{A}\in\R^{n\times n}$ is \emph{orthogonal} if
        $\mt{A}\T\mt{A}=\mt{A}\mt{A}\T=\mt{I}$. No inverse is mentioned.\see{D3}
  \item The converse ($\mt{A}\T\mt{A}=\mt{I}$ preserves length), and dot
        products and angles preserved too:
        $(\mt{A}\vc{x})\T(\mt{A}\vc{y})=\vc{x}\T\vc{y}$\prop.
\end{itemize}

\sect{21}{What an orthogonal matrix can do} In $\R^{2}$ the first column is a
unit vector $[\,a,c\,]$. The second is a unit vector orthogonal to it, and there
are exactly two: $\pm[\,-c,a\,]$. So every $2\times2$ orthogonal matrix is
$\begin{bmatrix}a&-c\\c&a\end{bmatrix}$ ($\det=+1$, a rotation) or
$\begin{bmatrix}a&c\\c&-a\end{bmatrix}$ ($\det=-1$, a reflection), with
$a^{2}+c^{2}=1$. Both satisfy $\mt{A}\mt{A}\T=\mt{I}$ by direct computation,
which settles D3 for $n=2$. Sort \S19's examples into the two families. The
projection is not orthogonal, because it destroys length. The angle form waits
for \S24.

\fits Catalog Problems~5--6 (Homework~3 F.2 gallery of motions, F.3 one motion
after another).

%------------------------------------------------------------------------------
\partplan{4}{Area, and the Rotation Matrix}
  {matrices\_4\_area\_and\_rotation}{\S22--\S25}{1--1.5}
%------------------------------------------------------------------------------
\sect{22}{Properties of the $2\times2$ determinant} Each is a one-line
computation: $\det\mt{I}=1$; $\det\mt{A}\T=\det\mt{A}$;
$\det(\alpha\mt{A})=\alpha^{2}\det\mt{A}$; swapping the columns (or rows) flips
the sign; linear in each column separately; adding a multiple of one column to
the other changes nothing; $\det(\mt{A}\mt{B})=\det\mt{A}\det\mt{B}$ (by
expansion); triangular and diagonal give the product of the diagonal;
orthogonal gives $\pm1$, since $\det(\mt{A}\T\mt{A})=(\det\mt{A})^{2}=1$, matching
\S21; $\det\neq0$ exactly at rank~2. Read each against the gallery: uniform
scaling $\alpha^{2}$, shear $1$, rotation $1$, reflection $-1$, projection $0$.

\sect{23}{The determinant is an area} Let $\vc{a}_1=[\,a,c\,]$ and
$\vc{a}_2=[\,b,d\,]$ be independent, with $\theta$ the angle between them
(vector unit \S10.3). The base is $\norm{\vc{a}_1}$ and the height is
$\norm{\vc{a}_2}\sin\theta$, so
$\text{Area}=\norm{\vc{a}_1}\norm{\vc{a}_2}\sin\theta$. Square it:
\[
  \text{Area}^{2}
  =\norm{\vc{a}_1}^{2}\norm{\vc{a}_2}^{2}(1-\cos^{2}\theta)
  =\norm{\vc{a}_1}^{2}\norm{\vc{a}_2}^{2}-(\vc{a}_1\cdot\vc{a}_2)^{2}
  =(a^{2}+c^{2})(b^{2}+d^{2})-(ab+cd)^{2}
  =(ad-bc)^{2},
\]
using $\cos\theta$ from \S11.3. Since $\sin\theta\ge0$ on $[0,\pi]$,
$\text{Area}=|ad-bc|=|\det\mt{A}|$. Consequences\prop:
\begin{itemize}[leftmargin=1.4em]
  \item The middle equality is Cauchy--Schwarz in $\R^{2}$ (vector unit \S13),
        with the gap measured exactly, so equality holds iff $\det=0$.
  \item $\text{Area}^{2}=\det(\mt{A}\T\mt{A})$, consistent with \S22.
  \item Every grid square scales by $|\det\mt{A}|$, so every region does.
  \item The product rule, seen geometrically: scale by $|\det\mt{B}|$, then by
        $|\det\mt{A}|$.
\end{itemize}

\sect{24}{The rotation matrix} Parameterize the unit circle as
$[\,\cos t,\sin t\,]$; rotation by $\theta$ sends angle $t$ to angle $t+\theta$.
So $\vc{e}_1$ ($t=0$) goes to $[\,\cos\theta,\sin\theta\,]$, and $\vc{e}_2$
($t=\pi/2$) goes to $[\,\cos(\theta+\frac{\pi}{2}),\sin(\theta+\frac{\pi}{2})\,]=[\,-\sin\theta,\cos\theta\,]$.
By \S13, $\mt{R}_\theta=\begin{bmatrix}\cos\theta&-\sin\theta\\\sin\theta&\cos\theta\end{bmatrix}$.
Then $\mt{R}_\theta\T\mt{R}_\theta=\mt{I}$ for all $\theta$, and
$\det\mt{R}_\theta=1$ for all $\theta$. This is \S21's
$\begin{bmatrix}a&-c\\c&a\end{bmatrix}$ family with $a=\cos\theta$,
$c=\sin\theta$. Also:
$\mt{R}_\alpha\mt{R}_\beta=\mt{R}_{\alpha+\beta}$ (``rotate by $\beta$, then
by $\alpha$'' is the product, by M1\see{D11}), which yields the
angle-addition formulas (or, given them, shows $\mt{R}_\theta$ turns every point of
the circle, not just $\vc{e}_1,\vc{e}_2$)\prop\see{D10}; $\mt{R}_\theta\T=\mt{R}_{-\theta}$; the
link to complex numbers \S6--\S7 (multiplying by $e^{i\theta}$)\prop; reflection
across the line at angle $\varphi$,
$\begin{bmatrix}\cos2\varphi&\sin2\varphi\\\sin2\varphi&-\cos2\varphi\end{bmatrix}$\prop.

\sect{25}{Rotations in $\R^{3}$} About the $z$-axis,
$\begin{bmatrix}\mt{R}_\theta&\vc{0}\\\vc{0}\T&1\end{bmatrix}$: it turns the
$xy$-plane and fixes the axis. About the $x$- and $y$-axes, with the sign
placement for $y$ explained. A rotation \emph{with respect to a plane} turns one
plane through the origin by $\theta$ and fixes the line orthogonal to it. In an
orthonormal basis adapted to the plane it is the block matrix above. Such a rotation is orthogonal,
with $\det=1$. Order matters: $\mt{R}_x\mt{R}_z\neq\mt{R}_z\mt{R}_x$, the
warning of \S8.4 made physical. Remark\prop: in $\R^{n}$, a rotation in the
$(i,j)$ coordinate plane (a Givens rotation).

\fits Catalog Problems~7--8 (Homework~3 F.4 rotations, F.5 shadows). Once
\S14 is lectured, F.2(b) no longer needs to avoid the word ``determinant''.

%------------------------------------------------------------------------------
\partplan{5}{Projection Matrices and the Direct Sum}
  {matrices\_5\_projections\_and\_direct\_sums}{\S26--\S29}{1--1.5}
%------------------------------------------------------------------------------
\sect{26}{The projection matrix onto a line} Start from the vector unit's
\S12.4, $\proj{\vc{u}}{\vc{v}}=\frac{\vc{u}\cdot\vc{v}}{\vc{u}\cdot\vc{u}}\vc{u}$.
Write the dot products as $\vc{u}\T\vc{v}$ (\S7.1) and move the scalar:
$\vc{u}\,\frac{\vc{u}\T\vc{v}}{\vc{u}\T\vc{u}}=\frac{\vc{u}\vc{u}\T}{\vc{u}\T\vc{u}}\vc{v}$.
So $\mt{P}_{\vc{u}}=\vc{u}\vc{u}\T/(\vc{u}\T\vc{u})$, an outer product (\S7.2)
over a number. Examples: $\vc{u}=\vc{e}_1$ gives \S19's $\diagop(1,0)$, and
$\vc{u}=[\,1,1\,]$ gives
$\frac12\begin{bmatrix}1&1\\1&1\end{bmatrix}$. Structure: symmetric, rank one,
$\colsp{P_{\vc{u}}}=\spanop\{\vc{u}\}$, $\det=0$,
$\mt{P}_{\vc{u}}\vc{u}=\vc{u}$, and $\mt{P}_{\vc{u}}\vc{w}=\vc{0}$ for
$\vc{w}\perp\vc{u}$. The same formula works unchanged in $\R^{n}$.

\sect{27}{Idempotence} Intuitively, doing it twice is the same as doing it once:
the shadow of a shadow is itself. $\mt{P}_{\vc{u}}^{2}=\mt{P}_{\vc{u}}$, since
$\vc{u}(\vc{u}\T\vc{u})\vc{u}\T/(\vc{u}\T\vc{u})^{2}=\mt{P}_{\vc{u}}$.
\textbf{Definition:} a square $\mt{P}$ is a \emph{projection matrix} if
$\mt{P}^{2}=\mt{P}$. Then\prop:
$\begin{bmatrix}1&1\\0&0\end{bmatrix}$ is idempotent and not symmetric. It
projects onto the $x$-axis \emph{along} $[\,1,-1\,]$, a slanted shadow. So an
\emph{orthogonal} projection matrix is one with $\mt{P}^{2}=\mt{P}$ and
$\mt{P}\T=\mt{P}$.

\sect{28}{The complementary projection} $\mt{I}-\mt{P}$ (not $\mt{Q}$; see N3):
$(\mt{I}-\mt{P})^{2}=\mt{I}-2\mt{P}+\mt{P}^{2}=\mt{I}-\mt{P}$;
$(\mt{I}-\mt{P}_{\vc{u}})\vc{v}=\vc{v}-\proj{\vc{u}}{\vc{v}}$, the orthogonal
piece of vector unit \S12.6; $\mt{P}(\mt{I}-\mt{P})=(\mt{I}-\mt{P})\mt{P}=\mt{0}$.
In $\R^{2}$, with $\vc{u}=[\,a,b\,]$ and $\vc{v}=[\,-b,a\,]\perp\vc{u}$,
$\vc{u}\vc{u}\T+\vc{v}\vc{v}\T=(a^{2}+b^{2})\mt{I}$, so
$\mt{I}-\mt{P}_{\vc{u}}=\mt{P}_{\vc{v}}$. The complement of the orthogonal
projection onto $\vc{u}$ is the orthogonal projection onto $\vc{v}$. Here
$\mt{P}_{\vc{v}}$ is defined from the projection of $\vc{u}$ onto $\vc{v}$,
$\mt{P}_{\vc{v}}\vc{u}=\proj{\vc{v}}{\vc{u}}$, as asked. Reflection across
$\spanop\{\vc{u}\}$ is $2\mt{P}_{\vc{u}}-\mt{I}=\mt{P}_{\vc{u}}-\mt{P}_{\vc{v}}$,
\S21's $\det=-1$ family again\prop.

\sect{29}{The direct sum} For orthogonal nonzero $\vc{u},\vc{v}\in\R^{2}$, every
$\vc{x}=\mt{P}_{\vc{u}}\vc{x}+\mt{P}_{\vc{v}}\vc{x}$, one piece on each line.
The split is unique because
$\spanop\{\vc{u}\}\cap\spanop\{\vc{v}\}=\{\vc{0}\}$. Intuitive definition:
$V=U\oplus W$ when every vector of $V$ is uniquely $\vc{u}+\vc{w}$ with
$\vc{u}\in U$, $\vc{w}\in W$; uniqueness amounts to $U\cap W=\{\vc{0}\}$. So
$\R^{2}=\spanop\{\vc{u}\}\oplus\spanop\{\vc{v}\}$, and
$\norm{\vc{x}}^{2}=\norm{\mt{P}_{\vc{u}}\vc{x}}^{2}+\norm{\mt{P}_{\vc{v}}\vc{x}}^{2}$
(vector unit \S23.4). Also\prop:
\begin{itemize}[leftmargin=1.4em]
  \item The oblique sum $\spanop\{[1,0]\}\oplus\spanop\{[1,1]\}$, split by
        \S27's idempotent: a direct sum need not be orthogonal.
  \item $\R^{3}$ as a line $\oplus$ a plane.
  \item Every projection gives $\colsp{P}\oplus\colsp{I-P}$.
\end{itemize}

\fits Catalog Problems~9--12 (Homework~3 Part~G). Rename $\mt{Q}_{\vc{u}}$
(N3) and rewrite G.3(e) as its instructor note says.

%------------------------------------------------------------------------------
\partplan{6}{The Null Space and Rank--Nullity}
  {matrices\_6\_null\_space\_and\_rank\_nullity}{\S30--\S34}{2}
%------------------------------------------------------------------------------
\sect{30}{Column and row space, revisited} Full rank, $\rank\mt{A}=\min(m,n)$;
full column rank (independent columns) and full row rank; rank-deficient
otherwise. $\mt{A}$ maps $\R^{n}\to\R^{m}$, landing in $\colsp{A}$;
$\mt{A}\T$ maps $\R^{m}\to\R^{n}$, landing in $\rowsp{A}$. The image of a
subspace is a subspace. For square $\mt{A}$, an \emph{invariant subspace} is a
$W$ with $\mt{A}\vc{w}\in W$ for all $\vc{w}\in W$. Examples:
\begin{itemize}[leftmargin=1.4em]
  \item $\spanop\{\vc{u}\}$ and $\spanop\{\vc{v}\}$ for $\mt{P}_{\vc{u}}$;
  \item the coordinate axes for a diagonal matrix;
  \item the $z$-axis and the $xy$-plane for a rotation about $z$;
  \item for $\mt{R}_\theta$ in $\R^{2}$ ($\theta\neq0,\pi$), only $\{\vc{0}\}$ and
        $\R^{2}$;
  \item $\colsp{A}$ (and, after \S31, $\nullsp{A}$) for any square $\mt{A}$.
\end{itemize}
Looking-ahead: a one-dimensional invariant subspace is spanned by an
eigenvector, and the eigen unit starts from it.

\sect{31}{The null space} $\nullsp{A}=\{\vc{x}:\mt{A}\vc{x}=\vc{0}\}$, the
\emph{kernel}; a subspace of $\R^{n}$; $\nullity(\mt{A})=\dim\nullsp{A}$.
$\mt{A}\vc{x}=\vc{0}$ is a homogeneous system (named here, solved in \S44).
Simple examples: $\mt{I}$; the zero matrix; $\mt{P}_{\vc{u}}$ (the line
$\spanop\{\vc{v}\}$); $\begin{bmatrix}1&2\\2&4\end{bmatrix}$
($\spanop\{[2,-1]\}$); a $2\times3$ (a line in $\R^{3}$); $\vc{u}\vc{v}\T$ (the
plane $\vc{v}^{\perp}$)\prop. Since
$\mt{A}\vc{x}=\mt{A}\vc{y}\iff\vc{x}-\vc{y}\in\nullsp{A}$, the kernel is what
$\mt{A}$ forgets. The solutions of $\mt{A}\vc{x}=\vc{b}$ are
$\vc{x}_p+\nullsp{A}$, which is the ``infinitely many'' outcome explained.
By the row reading, $\nullsp{A}\perp\rowsp{A}$: $\mt{A}\vc{x}=\vc{0}$ says
$\vc{x}$ is orthogonal to every row\prop.

\sect{32}{The rank--nullity theorem} For $\mt{A}\in\R^{m\times n}$,
$\rank\mt{A}+\nullity\mt{A}=n$. \emph{Proof:} take a basis
$\vc{k}_1,\dots,\vc{k}_p$ of $\nullsp{A}$ and extend it to a basis of $\R^{n}$
by $\vc{w}_1,\dots,\vc{w}_{n-p}$. Then $\mt{A}\vc{w}_1,\dots,\mt{A}\vc{w}_{n-p}$
is a basis of $\colsp{A}$. They span because $\mt{A}$ kills the $\vc{k}$'s. They
are independent because a dependence would put a combination of the $\vc{w}$'s in
$\nullsp{A}$.\see{D5} Many simple examples, in one table of matrix, $n$, rank,
nullity: each $2\times2$ rank, $\mt{P}_{\vc{u}}$, $\mt{R}_\theta$, a
$3\times3$ shear, $2\times3$, $3\times2$, $\vc{u}\vc{v}\T$, $\mt{0}$,
$\mt{I}$. Consequences:
\begin{itemize}[leftmargin=1.4em]
  \item For square $\mt{A}$: $\nullsp{A}=\{\vc{0}\}$ $\iff$ rank $n$ $\iff$
        $\colsp{A}=\R^{n}$. This is the fact Part~1 \S1.1 took on trust,
        modulo D5.
  \item A wide matrix always has nullity $\ge n-m>0$.
\end{itemize}

\sect{33}{A corollary: $\rank(\mt{A}\T\mt{A})=\rank(\mt{A})$} First,
$\nullsp{A^{\mathsf T}A}=\nullsp{A}$. One inclusion is clear. For the other,
$\mt{A}\T\mt{A}\vc{x}=\vc{0}$ gives
$\vc{x}\T\mt{A}\T\mt{A}\vc{x}=\norm{\mt{A}\vc{x}}^{2}=0$, so
$\mt{A}\vc{x}=\vc{0}$. Both matrices have $n$ columns, so rank--nullity gives
equal ranks. Second corollary\prop: $\rank\mt{A}\T=\rank\mt{A}$, i.e.\
$\dim\rowsp{A}=\dim\colsp{A}$. Since $\colsp{A^{\mathsf T}A}\subseteq\colsp{A\T}$,
$\rank\mt{A}\le\rank\mt{A}\T$; apply this to $\mt{A}\T$ for the reverse. No
elimination is needed, which settles D4. Third: $\mt{A}\T\mt{A}$ has full rank
$\iff$ $\mt{A}$ has full column rank, which is the hinge of \S49.

\sect{34}{The $2\times2$ inverse} Definition: $\mt{B}$ with
$\mt{A}\mt{B}=\mt{B}\mt{A}=\mt{I}$. Derive it by solving $\mt{A}\mt{X}=\mt{I}$
column by column:
$\mt{A}^{-1}=\frac{1}{ad-bc}\begin{bmatrix}d&-b\\-c&a\end{bmatrix}$. It exists
$\iff\det\neq0\iff$ rank~2 $\iff\nullsp{A}=\{\vc{0}\}$; one table ties these
four together. From the gallery:
\begin{itemize}[leftmargin=1.4em]
  \item $\mt{R}_\theta^{-1}=\mt{R}_{-\theta}=\mt{R}_\theta\T$;
  \item the inverse shear $\begin{bmatrix}1&-k\\0&1\end{bmatrix}$;
  \item $\diagop(1/s_1,1/s_2)$;
  \item a projection has none, because it forgets a line.
\end{itemize}
Solve \S1.5's system as $\vc{x}=\mt{V}^{-1}\vc{u}$, giving \S2.4's analogy with
$x=b/a$ its meaning.

%------------------------------------------------------------------------------
\partplan{7}{The Determinant in General}
  {matrices\_7\_the\_determinant}{\S35--\S39}{1.5}
%------------------------------------------------------------------------------
\sect{35}{Minors and cofactors} $M_{ij}$ is the determinant of $\mt{A}$ with row
$i$ and column $j$ deleted; $C_{ij}=(-1)^{i+j}M_{ij}$, the checkerboard of
signs. The $2\times2$ case as the base: $ad-bc=aC_{11}+bC_{12}$.

\sect{36}{Laplace (cofactor) expansion} Along row~1, recursively; then along any
row or column, all giving the same number (stated\see{D6}). Worked $3\times3$
examples; expand along the row or column with the most zeros. Also, as a
statement\prop: $|\det|$ of a $3\times3$ is the volume of the parallelepiped on
its columns, \S23 one dimension up.

\sect{37}{Principal minors} The determinant of the submatrix keeping the same
set of row and column indices; the leading principal minors (top-left
$k\times k$). Examples. Looking-ahead\prop: they come back in the eigen unit
(the characteristic polynomial) and in testing a Hessian (veccalc \S15).

\sect{38}{General properties} $\det\mt{I}=1$; $\det\mt{A}\T=\det\mt{A}$; a row
swap flips the sign; scaling one row by $\alpha$ scales $\det$ by $\alpha$, so
$\det(\alpha\mt{A})=\alpha^{n}\det\mt{A}$, which generalizes $\alpha^{2}$; adding a
multiple of one row to another changes nothing; two equal rows or a zero row
give $0$; $\det(\mt{A}\mt{B})=\det\mt{A}\det\mt{B}$; $\det\mt{A}\neq0\iff$
rank $n\iff\nullsp{A}=\{\vc{0}\}$. Which are proved and which are stated is D6.

\sect{39}{Special matrices, with proofs}
\begin{itemize}[leftmargin=1.4em]
  \item Triangular (expand down the first column; induction on $n$): the
        product of the diagonal. Diagonal is a special case.
  \item Orthogonal:
        $(\det\mt{Q})^{2}=\det\mt{Q}\T\det\mt{Q}=\det(\mt{Q}\T\mt{Q})=1$, so
        $\det\mt{Q}=\pm1$; \S25's rotations have $+1$.
  \item A projection $\mt{P}\neq\mt{I}$ has $\det=0$, since
        $\mt{P}(\mt{I}-\mt{P})=\mt{0}$ with $\mt{I}-\mt{P}\neq\mt{0}$.
\end{itemize}

%------------------------------------------------------------------------------
\partplan{8}{Solving Linear Systems}
  {matrices\_8\_gaussian\_elimination}{\S40--\S45}{2}
%------------------------------------------------------------------------------
\noindent\textit{Collects} Part~1 \S9's Looking-ahead (elimination, existence and
uniqueness) and \S4.4's hint that the method makes the matrix triangular.

\sect{40}{Elementary row operations} The augmented matrix
$[\,\mt{A}\mid\vc{b}\,]$; swap two rows, scale a row by a nonzero number, add a
multiple of one row to another. Each is reversible, so none changes the
solution set. Each is left multiplication by an elementary matrix (a
permutation, a scaling, a shear from Part~3)\prop.

\sect{41}{Triangular systems} Back substitution (upper triangular) and forward
substitution (lower), with pseudocode. They work iff no diagonal entry is zero,
i.e.\ $\det\neq0$ (\S39). Cost about $n^{2}$\prop.

\sect{42}{Gaussian elimination and row-echelon form} Pivots; a zero pivot forces
a row swap. Worked: Part~1 \S1.6's system
$\begin{bmatrix}1&2&0\\0&1&1\\1&0&1\end{bmatrix}\vc{x}=[\,5,5,4\,]$, solved at
last ($\vc{x}=[\,1,2,3\,]$). Then the $n\times n$ procedure.
$\det\mt{A}=\pm$ the product of the pivots, with the sign from the swaps, which
is far cheaper than \S36\prop.

\sect{43}{Gauss--Jordan elimination and reduced row-echelon form} Pivot and free
variables; every solution read off; a basis of $\nullsp{A}$ from the free
variables. Rank $=$ the number of pivots and nullity $=$ the number of free
variables. This is rank--nullity again, now by counting columns: the second
proof promised in D5.

\sect{44}{Homogeneous systems} $\mt{A}\vc{x}=\vc{0}$ always has $\vc{0}$; it
has nontrivial solutions $\iff$ there is a free variable $\iff$ nullity $>0$,
which is guaranteed when there are more unknowns than equations. General solution =
particular + homogeneous (\S31). Consistency:
$\vc{b}\in\colsp{A}\iff\rank[\,\mt{A}\mid\vc{b}\,]=\rank\mt{A}$. The three
outcomes, settled.

\sect{45}{Overdetermined and underdetermined systems} With $m>n$ (more equations
than unknowns), a system is usually inconsistent, which leads to least squares in
Part~9. With $m<n$, a consistent system has infinitely many solutions (nullity
$\ge n-m$). Which one to pick is the minimum-norm question of \S54.

%------------------------------------------------------------------------------
\partplan{9}{Least Squares}
  {matrices\_9\_least\_squares}{\S46--\S49}{1.5}
%------------------------------------------------------------------------------
\sect{46}{The closest point on a line} For vectors $\vc{u},\vc{v}\in\R^{n}$ and a
scalar $x$, minimize $f(x)=\norm{\vc{v}-x\vc{u}}_2^{2}
=\vc{v}\T\vc{v}-2x\,\vc{u}\T\vc{v}+x^{2}\,\vc{u}\T\vc{u}$, a parabola in $x$.
$f'(x)=0$ gives $\hat x=\vc{u}\T\vc{v}/\vc{u}\T\vc{u}$, the coefficient of
$\proj{\vc{u}}{\vc{v}}$ (vector unit \S12.4). The residual is the orthogonal
piece (\S12.6), and the Pythagorean identity shows every other $x$ does worse.
This needs one-variable calculus only (veccalc \S1).

\sect{47}{The closest point in a column space} Let $\vc{b}\notin\colsp{A}$
(\S45) and minimize $\norm{\mt{A}\vc{x}-\vc{b}}^{2}$. As $\vc{x}$ varies,
$\mt{A}\vc{x}$ sweeps $\colsp{A}$. The closest point
$\mt{A}\hat{\vc{x}}$ is where the residual is orthogonal to $\colsp{A}$, i.e.\
to every column: $\vc{a}_j\T(\vc{b}-\mt{A}\hat{\vc{x}})=0$ for all $j$. That
gives the \textbf{normal equations}
$\mt{A}\T\mt{A}\hat{\vc{x}}=\mt{A}\T\vc{b}$. One column recovers \S46.

\sect{48}{The same answer by calculus} Expand
$f(\vc{x})=\vc{x}\T\mt{A}\T\mt{A}\vc{x}-2\vc{b}\T\mt{A}\vc{x}+\vc{b}\T\vc{b}$.
Componentwise, $\nabla(\vc{c}\T\vc{x})=\vc{c}$ and
$\nabla(\vc{x}\T\mt{M}\vc{x})=2\mt{M}\vc{x}$ for symmetric $\mt{M}$. So
$\nabla f=2\mt{A}\T\mt{A}\vc{x}-2\mt{A}\T\vc{b}=\vc{0}$ is the normal equations.
The Hessian is $2\mt{A}\T\mt{A}$, and
$\vc{x}\T\mt{A}\T\mt{A}\vc{x}=\norm{\mt{A}\vc{x}}^{2}\ge0$, so the critical point
is a minimum.\see{D9}

\sect{49}{Solving the normal equations; fitting a line} The normal equations are
always consistent\prop:
$\mt{A}\T\vc{b}\in\colsp{A\T}=\colsp{A^{\mathsf T}A}$. The two spaces are nested
and have the same dimension, by \S33. The solution is unique $\iff$ $\mt{A}$ has
full column rank (\S33). Solve them by elimination (Part~8); no inverse is used.\see{D7}
Worked example: the line $y=c_0+c_1t$ through three or four data points (linear
regression), with the residual checked orthogonal to both columns.

\fits Catalog Problems~13--14 (Homework~3 Part~H, the Hessian). H.2, the
Hessian of a quadratic function, is exactly \S48's computation.

%------------------------------------------------------------------------------
\partplan{10}{Cramer's Rule, the Inverse, and the Pseudo-inverse}
  {matrices\_10\_the\_inverse\_and\_the\_pseudoinverse}{\S50--\S54}{2}
%------------------------------------------------------------------------------
\sect{50}{Cramer's rule} For $\det\mt{A}\neq0$,
$x_i=\det\mt{A}_i(\vc{b})/\det\mt{A}$, where $\mt{A}_i(\vc{b})$ has column $i$
replaced by $\vc{b}$. Geometric intuition ($2\times2$): since
$\vc{b}=x_1\vc{a}_1+x_2\vc{a}_2$,
$\det[\,\vc{b}\ \vc{a}_2\,]=x_1\det[\,\vc{a}_1\ \vc{a}_2\,]+x_2\det[\,\vc{a}_2\
\vc{a}_2\,]=x_1\det\mt{A}$. Sliding $\vc{b}$ parallel to $\vc{a}_2$ is a shear
and keeps the area (\S22--\S23). So the parallelogram on $(\vc{b},\vc{a}_2)$ has
$x_1$ times the area of the one on $(\vc{a}_1,\vc{a}_2)$. For $n\times n$, the
same argument works with volumes. Cost remark\prop: $n+1$ determinants against
one elimination.

\sect{51}{The inverse of an $n\times n$ matrix} Definition and uniqueness. The
following are equivalent: $\mt{A}$ is invertible; $\det\neq0$; rank $n$;
$\nullsp{A}=\{\vc{0}\}$; $\mt{A}\vc{x}=\vc{b}$ has exactly one solution for every
$\vc{b}$ (collecting Parts~2, 6, 7, 8). For square matrices, a one-sided
inverse is enough.
\begin{itemize}[leftmargin=1.4em]
  \item $\mt{B}\mt{A}=\mt{I}$ gives $\nullsp{A}=\{\vc{0}\}$, so rank $n$.
  \item Then $\mt{A}\mt{C}=\mt{I}$ is solvable column by column.
  \item So $\mt{B}=\mt{B}(\mt{A}\mt{C})=(\mt{B}\mt{A})\mt{C}=\mt{C}$.
\end{itemize}
This settles D3: $\mt{Q}\T\mt{Q}=\mt{I}$ alone forces $\mt{Q}\mt{Q}\T=\mt{I}$,
and $\mt{Q}^{-1}=\mt{Q}\T$. Computing $\mt{A}^{-1}$ by Gauss--Jordan on
$[\,\mt{A}\mid\mt{I}\,]$.

\sect{52}{The inverse by the adjugate} $\adj\mt{A}=\mt{C}\T$, the cofactor
matrix transposed, and $\mt{A}\,\adj\mt{A}=\det(\mt{A})\,\mt{I}$. On the diagonal
this is Laplace expansion; off the diagonal it is the expansion of a matrix with
two equal rows, which is $0$. So $\mt{A}^{-1}=\adj\mt{A}/\det\mt{A}$, which is
Cramer's rule applied to each column of $\mt{I}$. \S34's formula is the
$2\times2$ case; a $3\times3$ example.\see{N4}

\sect{53}{Properties of the inverse} $(\mt{A}^{-1})^{-1}=\mt{A}$;
$(\mt{A}\mt{B})^{-1}=\mt{B}^{-1}\mt{A}^{-1}$ (order reverses, like M6);
$(\mt{A}\T)^{-1}=(\mt{A}^{-1})\T$; $(\alpha\mt{A})^{-1}=\alpha^{-1}\mt{A}^{-1}$;
$\det\mt{A}^{-1}=1/\det\mt{A}$; diagonal matrices invert entrywise, and
triangular ones stay triangular; $\mt{Q}^{-1}=\mt{Q}\T$;
$\mt{R}_\theta^{-1}=\mt{R}_{-\theta}$. To solve a system, eliminate rather than
forming $\mt{A}^{-1}$\prop. With full column rank,
$\hat{\vc{x}}=(\mt{A}\T\mt{A})^{-1}\mt{A}\T\vc{b}$.

\sect{54}{The pseudo-inverse} The cases:
\begin{itemize}[leftmargin=1.4em]
  \item \textbf{Full column rank:} $\mt{A}^{+}=(\mt{A}\T\mt{A})^{-1}\mt{A}\T$,
        from \S49, with $\mt{A}^{+}\mt{A}=\mt{I}_n$, and
        $\mt{A}\mt{A}^{+}=\mt{A}(\mt{A}\T\mt{A})^{-1}\mt{A}\T$ is the orthogonal
        projection onto $\colsp{A}$. For a single column $\mt{A}=\vc{u}$,
        $\vc{u}^{+}=\vc{u}\T/\vc{u}\T\vc{u}$ and $\vc{u}\vc{u}^{+}=\mt{P}_{\vc{u}}$
        of \S26.
  \item \textbf{Full row rank:} $\mt{A}^{+}=\mt{A}\T(\mt{A}\mt{A}\T)^{-1}$
        gives the minimum-norm solution of \S45's underdetermined system. It is
        the solution in $\rowsp{A}$, orthogonal to $\nullsp{A}$.
  \item \textbf{Invertible:} $\mt{A}^{+}=\mt{A}^{-1}$.
\end{itemize}
\textbf{Moore--Penrose:} $\mt{A}^{+}$ is the unique $\mt{X}$ with
$\mt{A}\mt{X}\mt{A}=\mt{A}$, $\mt{X}\mt{A}\mt{X}=\mt{X}$,
$(\mt{A}\mt{X})\T=\mt{A}\mt{X}$, $(\mt{X}\mt{A})\T=\mt{X}\mt{A}$. Check both
formulas against the four conditions. The rank-deficient case is stated here
and computed with the SVD in L3\see{D8}. One rank-one example\prop:
$(\vc{x}\vc{y}\T)^{+}=\vc{y}\vc{x}\T/(\norm{\vc{x}}^{2}\norm{\vc{y}}^{2})$.

%------------------------------------------------------------------------------
\partplan{11}{Gram--Schmidt and the $QR$ Decomposition}
  {matrices\_11\_gram\_schmidt\_and\_qr}{\S55--\S58}{1.5}
%------------------------------------------------------------------------------
\noindent\textit{Collects} computational unit \S20.3's closing promise: ``a
factorisation of the matrix whose columns are the $\vc{v}_k$ \ldots\ We will
return to Gram--Schmidt.''

\sect{55}{Gram--Schmidt in matrix language} \see{D9} The algorithm is derived in
full in computational unit Part~4 (\S16--\S19). Here it is restated in matrix
terms: $\vc{w}_k=\vc{a}_k-\sum_{j<k}(\vc{q}_j\T\vc{a}_k)\vc{q}_j$ and
$\vc{q}_k=\vc{w}_k/\norm{\vc{w}_k}$. It requires independent columns
(computational \S19).

\sect{56}{The $QR$ decomposition} Rearranged,
$\vc{a}_k=\sum_{j\le k}r_{jk}\vc{q}_j$ with $r_{jk}=\vc{q}_j\T\vc{a}_k$ and
$r_{kk}=\norm{\vc{w}_k}>0$. So $\mt{A}=\mt{Q}\mt{R}$, where
$\mt{Q}\in\R^{m\times n}$ has orthonormal columns ($\mt{Q}\T\mt{Q}=\mt{I}_n$) and
$\mt{R}\in\R^{n\times n}$ is upper triangular with a positive diagonal. Worked
example: computational \S18.3's $[1,1,0],[1,0,1],[0,1,1]$, whose $\vc{q}$'s are
already computed there, with $\mt{R}=\mt{Q}\T\mt{A}$. For square $\mt{A}$,
$\mt{Q}$ is orthogonal (Part~3). For tall $\mt{A}$, $\mt{Q}\T\mt{Q}=\mt{I}$ but
$\mt{Q}\mt{Q}\T\neq\mt{I}$: it is the projection onto $\colsp{A}$ (\S54).

\sect{57}{$\colsp{Q}=\colsp{A}$} $\mt{A}=\mt{Q}\mt{R}$ gives
$\colsp{A}\subseteq\colsp{Q}$. $\mt{R}$ is triangular with a nonzero diagonal, so
$\det\mt{R}\neq0$ (\S39) and $\mt{Q}=\mt{A}\mt{R}^{-1}$, giving the reverse
inclusion. More strongly,
$\spanop\{\vc{q}_1,\dots,\vc{q}_k\}=\spanop\{\vc{a}_1,\dots,\vc{a}_k\}$ for every
$k$.

\sect{58}{$QR$ and least squares}
$\mt{A}\T\mt{A}=\mt{R}\T\mt{Q}\T\mt{Q}\mt{R}=\mt{R}\T\mt{R}$, so the normal
equations reduce to $\mt{R}\hat{\vc{x}}=\mt{Q}\T\vc{b}$, a single back
substitution (\S41). Also $\mt{A}^{+}=\mt{R}^{-1}\mt{Q}\T$, and the projection
onto $\colsp{A}$ is $\mt{Q}\mt{Q}\T$. Redo \S49's line fit with $QR$. Remark\prop:
why $QR$ is preferred in floating point (forming $\mt{A}\T\mt{A}$ squares the
sensitivity), and classical versus modified Gram--Schmidt (computational
\S20.3).

%==============================================================================
\section{The eigen unit (the final)}
%==============================================================================
This is a new unit, so section numbering restarts at \S1, the convention matrix
Part~1 set for a new unit. Every topic on the list is placed below. Several
matrix-unit sections plant it on purpose: \S30 (invariant subspaces), \S37
(principal minors), \S38--\S39 ($\det(\mt{A}-\lambda\mt{I})$, triangular
determinants), \S26--\S29 (projections, direct sums), \S24 ($\mt{R}_\theta$ has
no real eigenvector), \S53 ($\mt{V}^{-1}$), and Part~1 \S9.3
($\tr\mt{A}\mt{B}=\tr\mt{B}\mt{A}$).

\partplan{E1}{Eigenvalues, Eigenvectors, and the Characteristic Polynomial}
  {eigen\_1\_eigenvalues\_and\_the\_characteristic\_polynomial}{\S1--\S4}{1.5}
\sect{1}{Eigenpairs} $\mt{A}\vc{v}=\lambda\vc{v}$, $\vc{v}\neq\vc{0}$; an
eigenvector spans a one-dimensional invariant subspace (matrix \S30). The
gallery first, by geometry:
\begin{itemize}[leftmargin=1.4em]
  \item scalings;
  \item projections ($1$ on the column space, $0$ on the null space);
  \item reflections ($\pm1$);
  \item shears (only $1$);
  \item $\mt{R}_\theta$ (no real eigenvalue).
\end{itemize}

\sect{2}{The characteristic polynomial} $\lambda$ is an eigenvalue
$\iff\nullsp{A-\lambda I}\neq\{\vc{0}\}\iff\det(\mt{A}-\lambda\mt{I})=0$
(matrix \S38). The eigenspace is $\nullsp{A-\lambda I}$, and the polynomial has degree $n$.

\sect{3}{The $2\times2$ case; trace and determinant} With $\tau=\tr\mt{A}$ and $\delta=\det\mt{A}$,
$p(\lambda)=\lambda^{2}-\tau\lambda+\delta$, so
$\lambda=\tfrac12\bigl(\tau\pm\sqrt{\tau^{2}-4\delta}\bigr)$. The
discriminant decides between two real roots, a repeated root, or a complex
pair. $\mt{R}_\theta$ gives $e^{\pm i\theta}$ (complex numbers \S6).
\textbf{Theorem:} $\tr\mt{A}=\sum\lambda_i$ and $\det\mt{A}=\prod\lambda_i$
(complex, with multiplicity). Proofs:
\begin{itemize}[leftmargin=1.4em]
  \item $2\times2$: the sum and product of the two roots above.
  \item $n\times n$: $\det(\lambda\mt{I}-\mt{A})=\prod(\lambda-\lambda_i)$ (the
        fundamental theorem of algebra, stated). Setting $\lambda=0$ gives the
        determinant. The $\lambda^{n-1}$ coefficient gives the trace, because
        only the diagonal product contributes: in a first-row expansion, every
        other term's minor lacks two diagonal entries.
  \item For diagonalizable $\mt{A}$, also
        $\tr(\mt{V}\Lambda\mt{V}^{-1})=\tr(\Lambda)$ by R4\prop.
\end{itemize}

\sect{4}{Abel--Ruffini} Statement: no formula in radicals gives the roots of a
general polynomial of degree $\ge5$. The relationship: every monic polynomial
is the characteristic polynomial of its companion matrix\prop. So a radical
formula for the eigenvalues of every $5\times5$ matrix would solve every
quintic. Eigenvalues from $5\times5$ up must therefore, in general, be computed
by iteration rather than formula (looking ahead to the power method and the $QR$
algorithm).

\partplan{E2}{Diagonalization and the Spectral Theorem}
  {eigen\_2\_diagonalization\_and\_the\_spectral\_theorem}{\S5--\S7}{1.5}
\sect{5}{Diagonalization} $\mt{A}=\mt{V}\Lambda\mt{V}^{-1}$, with eigenvectors as the
columns of $\mt{V}$. This holds $\iff$ there are $n$ independent eigenvectors.
Distinct eigenvalues give independent eigenvectors. The shear
$\begin{bmatrix}1&1\\0&1\end{bmatrix}$ is not diagonalizable. Powers:
$\mt{A}^{k}=\mt{V}\Lambda^{k}\mt{V}^{-1}$. Presented algebraically. Reading it
as a change of basis, ``in the eigenbasis $\mt{A}$ acts as $\Lambda$'', is
planted here and done in L2.\see{D12}

\sect{6}{The spectral theorem} For real symmetric $\mt{A}$:
\begin{itemize}[leftmargin=1.4em]
  \item the eigenvalues are real;
  \item eigenvectors of distinct eigenvalues are orthogonal (short proofs of
        both);
  \item there is an orthonormal eigenbasis, with
        $\mt{A}=\mt{Q}\Lambda\mt{Q}\T=\sum\lambda_i\vc{q}_i\vc{q}_i\T$, a
        weighted sum of the rank-one projections of matrix \S26;
  \item $\R^{n}$ is the direct sum of the eigenspaces (matrix \S29).
\end{itemize}
Existence of the eigenbasis in general is stated. Applied to the symmetric
matrix $\mt{A}\T\mt{A}$, this is the engine of the SVD (L3).

\sect{7}{The spectral radius} $\rho(\mt{A})=\max|\lambda_i|$. For diagonalizable
$\mt{A}$, $\mt{A}^{k}\to\mt{0}\iff\rho(\mt{A})<1$\prop.

\partplan{E3}{Structure and Location}
  {eigen\_3\_structure\_and\_gershgorin}{\S8--\S9}{1}
\sect{8}{Structure and eigenvalues} A table for real $n\times n$ matrices, each
row with a short proof:
\begin{itemize}[leftmargin=1.4em]
  \item triangular and diagonal: the diagonal entries (matrix \S39 applied to
        $\mt{A}-\lambda\mt{I}$);
  \item symmetric: real eigenvalues, orthogonal eigenvectors;
  \item skew-symmetric: $0$ or purely imaginary;
  \item orthogonal: $|\lambda|=1$;
  \item projection: $0$ and $1$; reflection: $\pm1$; nilpotent: only $0$;
  \item $\vc{u}\vc{v}\T$: $\vc{v}\T\vc{u}$ and zeros;
  \item $\mt{A}\T$: the same eigenvalues;
  \item $\mt{A}^{k}$, $\mt{A}^{-1}$, $\mt{A}+c\mt{I}$: $\lambda^{k}$,
        $1/\lambda$, $\lambda+c$;
  \item $\mt{A}\T\mt{A}$: eigenvalues $\ge0$;
  \item invertible $\iff$ $0$ is not an eigenvalue.
\end{itemize}

\sect{9}{The Gershgorin circle theorem} Every eigenvalue lies in a disc centered
at some $a_{ii}$ with radius $R_i=\sum_{j\neq i}|a_{ij}|$ (discs in
$\mathbb{C}$: complex numbers \S5.3). \emph{Proof:} in
$\mt{A}\vc{v}=\lambda\vc{v}$, choose $i$ with $|v_i|$ largest. Examples.
Corollaries\prop: a strictly diagonally dominant matrix is invertible, and
$\rho(\mt{A})\le\max_i\sum_j|a_{ij}|$.

%==============================================================================
\section{The linear-maps unit (the end of the course)}
%==============================================================================
This is a new unit, restarting at \S1. Much of it reinterprets what is already
built, so it is shorter than its importance. It collects three debts:
\begin{itemize}[leftmargin=1.4em]
  \item Part~1 \S8's ``no attempt is made yet to justify why the rule is this
        one; that motivation comes later in the course'';
  \item D8's deferred pseudo-inverse;
  \item the change-of-basis reading of diagonalization planted in E2.
\end{itemize}

\partplan{L1}{Linear Maps and Composition}
  {linmaps\_1\_linear\_maps\_and\_composition}{\S1--\S3}{1}
\sect{1}{Linear maps} $T:\R^{n}\to\R^{m}$ with $T(\vc{x}+\vc{y})=T\vc{x}+T\vc{y}$
and $T(\alpha\vc{x})=\alpha T\vc{x}$. Examples and non-examples, including
translation (matrix \S17). Kernel and image of a map, which are matrix
\S30--\S32 restated for maps: injective $\iff$ kernel $\{\vc{0}\}$, surjective
$\iff$ image $=\R^{m}$, invertible $\iff$ both.

\sect{2}{The matrix of a linear map} $[T]=[\,T\vc{e}_1\mid\cdots\mid T\vc{e}_n\,]$,
matrix \S13 made general. Every linear map $\R^{n}\to\R^{m}$ is multiplication
by exactly one matrix, and vice versa.

\sect{3}{Composition, and where matrix multiplication comes from} Let
$T:\R^{n}\to\R^{p}$ and $S:\R^{p}\to\R^{m}$. Then $S\circ T$ is linear, and
column $j$ of its matrix is $S(T\vc{e}_j)=[S]\,(\text{column } j\text{ of }[T])$.
Writing that out entry by entry \emph{forces}
$c_{ij}=\sum_k a_{ik}b_{kj}$, the rule of Part~1 \S8.1, now derived rather than
decreed. The shape rule $(m\times p)(p\times n)$ is the requirement that the
maps chain. M1 is associativity of composing functions. Non-commutativity is
order of operations: rotate-then-reflect differs from reflect-then-rotate
(Part~3). $(\mt{A}\mt{B})^{-1}=\mt{B}^{-1}\mt{A}^{-1}$ is undoing in reverse
order.

\partplan{L2}{Change of Basis}
  {linmaps\_2\_change\_of\_basis}{\S4--\S5}{1}
\sect{4}{Coordinates in a basis} This returns to matrix Part~1 \S1--\S2:
$\vc{u}=\mt{V}\vc{x}$, so $\vc{x}=\mt{V}^{-1}\vc{u}$ (matrix \S51). Part~1's
distinction between the point $\vc{u}$ and its recipe $\vc{x}$ becomes the
change-of-basis matrix. Changing between two non-standard bases:
$\mt{W}^{-1}\mt{V}$. For an orthonormal basis, $\mt{V}^{-1}=\mt{V}\T$, so
coordinates are dot products (computational \S16).

\sect{5}{The matrix of a map in another basis; similarity} In the basis of
$\mt{V}$'s columns, the map $\vc{x}\mapsto\mt{A}\vc{x}$ has matrix
$\mt{V}^{-1}\mt{A}\mt{V}$. $\mt{A}$ and $\mt{B}$ are \emph{similar} when
$\mt{B}=\mt{V}^{-1}\mt{A}\mt{V}$. Similar matrices share determinant, trace,
rank, and eigenvalues, each a one-line proof (the product rule, R4,
$\det(\mt{V}^{-1}(\mt{A}-\lambda\mt{I})\mt{V})$). Diagonalization reread: in
the eigenbasis, $\mt{A}$ \emph{is} $\Lambda$. The spectral theorem reread:
$\mt{Q}\T\mt{A}\mt{Q}=\Lambda$, a change to an orthonormal basis.

\partplan{L3}{The Singular Value Decomposition}
  {linmaps\_3\_the\_svd}{\S6--\S8}{1.5}
\sect{6}{The circle goes to an ellipse} Any $2\times2$ matrix sends the unit
circle to an ellipse (possibly flattened). It looks like non-uniform scaling
(matrix \S16) between two rotations or reflections (\S21). The claim is that
every matrix is exactly that: \emph{orthogonal, scale, orthogonal}.

\sect{7}{Construction} $\mt{A}\T\mt{A}$ is symmetric with eigenvalues $\ge0$
(E2, E3). Take orthonormal eigenvectors $\vc{v}_i$, set
$\sigma_i=\sqrt{\lambda_i}$, and for $\sigma_i>0$ set
$\vc{u}_i=\mt{A}\vc{v}_i/\sigma_i$. These are orthonormal:
$\vc{u}_i\T\vc{u}_j=\vc{v}_i\T\mt{A}\T\mt{A}\vc{v}_j/(\sigma_i\sigma_j)=\delta_{ij}$.
So $\mt{A}=\mt{U}\Sigma\mt{V}\T$. The number of nonzero $\sigma_i$ is
$\rank(\mt{A}\T\mt{A})=\rank\mt{A}$ (matrix \S33). Thin and full forms. A
$2\times2$ example, worked by hand.

\sect{8}{What it gives}
\begin{itemize}[leftmargin=1.4em]
  \item Orthonormal bases of all four subspaces. The first $r$ $\vc{u}$'s span
        $\colsp{A}$. The first $r$ $\vc{v}$'s span $\rowsp{A}$. The rest span
        $\nullsp{A}$ and $\nullsp{A\T}$. This is matrix \S30--\S33 made
        explicit.
  \item $\mt{A}=\sum\sigma_i\vc{u}_i\vc{v}_i\T$, a sum of rank-one matrices.
  \item The pseudo-inverse in every case,
        $\mt{A}^{+}=\mt{V}\Sigma^{+}\mt{U}\T$. This collects D8; check it
        against the four conditions of \S54.
  \item For square $\mt{A}$, $|\det\mt{A}|=\prod\sigma_i$: area again, \S23\prop.
  \item Truncating the sum gives the best low-rank approximation (Eckart--Young,
        stated only; it is listed among the not-guaranteed topics)\prop.
\end{itemize}

%==============================================================================
\section{Decisions to settle before writing}
%==============================================================================
\begin{enumerate}[label=\textbf{D\arabic*.},leftmargin=2.4em,itemsep=5pt]
  \item \textbf{Part~1 says elimination is ``next''.} Its closing Looking-ahead
        (end of \S9) promises Gaussian elimination as the next topic; under
        this order it is Part~8. Recommended: change ``it is next'' to ``it is
        coming'' in Part~1 (a one-line edit and recompile of a distributed
        document), and have Part~2 open by explaining the detour: when a
        solution exists comes before how to find it.
  \item \textbf{Translation is not linear.} No matrix acting on $\R^{2}$
        translates, because $\mt{A}\vc{0}=\vc{0}$ always. \S17 proves this, then
        translates with a $3\times3$ matrix in homogeneous coordinates.
        Alternatives: keep translation only as the non-example, or drop it.
  \item \textbf{$\mt{A}\T\mt{A}=\mt{I}$ is all the length argument gives.} The
        definition as specified ($\mt{A}\T\mt{A}=\mt{A}\mt{A}\T=\mt{I}$) is
        kept. For $2\times2$, the second identity follows from the first by
        \S21's classification. For $n\times n$ it follows from ``a one-sided
        inverse of a square matrix is two-sided'', proved in \S51 from
        rank--nullity. Part~3 says so in a Looking-ahead.
  \item \textbf{Row rank $=$ column rank.} Before \S33 this is checked for
        $2\times2$ only: $ad-bc$ is the same number read along the rows. \S33
        proves it in general as a corollary of
        $\rank(\mt{A}\T\mt{A})=\rank\mt{A}$, which keeps the given order and
        needs no elimination.
  \item \textbf{Rank--nullity before elimination.} The \S32 proof extends a
        basis, so it takes as given three facts about dimension that no
        document proves in general: every subspace of $\R^{n}$ has a basis; an
        independent set extends to a basis; all bases of a subspace have the
        same size. This is the same footing on which Part~1 \S1.1 accepted
        ``$n$ independent vectors span $\R^{n}$''. \S43 gives the second,
        computational proof by counting pivots.
  \item \textbf{What is proved about the general determinant.} Stated without
        proof (the proofs need permutations): every row or column expansion
        gives the same number; $\det\mt{A}\T=\det\mt{A}$; the $n\times n$
        product rule. Proved: $\det(\alpha\mt{A})=\alpha^{n}\det\mt{A}$; a zero
        row; two equal rows (from the swap rule); triangular and diagonal (by
        induction); orthogonal (from the product and transpose rules);
        $\mt{A}\,\adj\mt{A}=\det(\mt{A})\,\mt{I}$.
  \item \textbf{Least squares comes before the general inverse.} So Part~9
        solves the normal equations by elimination and never writes
        $(\mt{A}\T\mt{A})^{-1}$. That formula arrives in \S53 and becomes
        $\mt{A}^{+}$ in \S54.
  \item \textbf{Scope of the pseudo-inverse.} For Midterm~2: full column rank,
        full row rank, invertible, and the four Moore--Penrose conditions.
        Existence and uniqueness for rank-deficient $\mt{A}$ is stated in \S54
        and computed in L3 as $\mt{V}\Sigma^{+}\mt{U}\T$, so it is on the final.
  \item \textbf{Prerequisites not yet lectured.} As of the last recorded lecture
        (Oct~8), these have not been covered:
        \begin{itemize}[leftmargin=1.2em]
          \item veccalc Part~2 (\S10--\S15: the gradient in $\R^{n}$ and the
                Hessian), needed by \S48;
          \item computational Part~4 (Gram--Schmidt; the Kronecker delta of
                \S16.2), needed by \S20 and Part~11;
          \item complex numbers \S5.3 and \S6 (discs, Euler's formula), used by
                the eigen unit.
        \end{itemize}
        Recommended: restate $\delta_{ij}$ in \S20 and the two gradient facts in
        \S48 in look-back boxes, as Homework~3's Hessian part did. Lecture
        computational Part~4's \S17--\S18 instead of rewriting Gram--Schmidt.
  \item \textbf{Angle-addition formulas} appear in no course document. \S24
        states them, or derives them from
        $\mt{R}_\alpha\mt{R}_\beta=\mt{R}_{\alpha+\beta}$.
  \item \textbf{Composition is used before it is derived.} The matrix unit
        reads $\mt{A}\mt{B}$ as ``do $\mt{B}$, then $\mt{A}$'' in \S24
        ($\mt{R}_\alpha\mt{R}_\beta$), \S25 ($\mt{R}_x\mt{R}_z$), and the
        Part~3 compositions. That reading is free from M1:
        $(\mt{A}\mt{B})\vc{x}=\mt{A}(\mt{B}\vc{x})$. L1 reverses the logic and
        derives the product rule \emph{from} composition, as specified. An
        alternative is to move L1~\S3 forward to sit beside \S13, where the
        gallery first composes motions. That would answer Part~1 \S8's
        ``why this rule'' as soon as it is needed and shorten the end of term,
        but it departs from the given order.
  \item \textbf{Change of basis comes after diagonalization.} In this order,
        E2 presents $\mt{V}\Lambda\mt{V}^{-1}$ algebraically, and L2 supplies
        the change-of-basis meaning afterwards. This works because Part~1
        \S1--\S2 already treated $\mt{V}\vc{x}=\vc{u}$ as ``coordinates in a
        basis''. E2 only needs a Looking-ahead.
\end{enumerate}

%==============================================================================
\section{Notation}
%==============================================================================
\begin{enumerate}[label=\textbf{N\arabic*.},leftmargin=2.4em,itemsep=4pt]
  \item $\colsp{A}$ and $\rowsp{A}$, as used in class. Proposed: $\nullsp{A}$
        in the same style, with ``kernel'' and the textbook spellings
        $N(\mt{A})$, $\ker\mt{A}$ mentioned once. Subspaces take brackets and
        numbers take parentheses: $\rank(\mt{A})$, $\nullity(\mt{A})$,
        $\det(\mt{A})$, $\tr(\mt{A})$.
  \item Never $|\mt{A}|$ for the determinant: it collides with
        $|\det\mt{A}|$ in \S23.
  \item $\mt{Q}$ is reserved for orthogonal matrices and $QR$. The
        complementary projection is written $\mt{I}-\mt{P}_{\vc{u}}$. The
        held-back Homework~3 Part~G problems (catalog 9--12) write
        $\mt{Q}_{\vc{u}}$ for it; rename it when they are assigned.
  \item ``Adjoint'': use \emph{adjugate}, $\adj\mt{A}$, with one sentence
        noting that ``(classical) adjoint'' is the older name and that
        ``adjoint'' elsewhere usually means the (conjugate) transpose.
  \item $\mt{R}_\theta$ for rotations, and $\mt{R}$ (no subscript) for the triangular
        factor of $QR$. $\mt{P}_{\vc{u}}$ for the orthogonal projection onto
        $\spanop\{\vc{u}\}$.
  \item Part~2 opens with \verb|\setcounter{section}{9}|, as the vector and
        computational units do. The eigen and linear-maps units each restart
        at \S1.
\end{enumerate}

%==============================================================================
\section{Pacing against the calendar}\label{sec:pacing}
%==============================================================================
Meetings left after Oct~22 (whatever Oct~20 and~22 covered is not recorded):
Oct~27, 29; Nov~3, 5, 10, 12, 17, 19, 24; Dec~1, 3, 8, 10. That is
\textbf{13 meetings}.

\medskip
\noindent\textbf{Estimated demand}, at the pace so far (matrix Part~1 took two
meetings; vector-unit Parts took 1.5 to 2.5):
\begin{itemize}[leftmargin=1.4em]
  \item matrix Parts~2--11: about 16 meetings;
  \item the eigen unit: about 4;
  \item the linear-maps unit: about 3.5;
  \item Midterm~2: 1.
\end{itemize}
That is about \textbf{24 meetings} against 13. Even at a full Part every
meeting, the sixteen Parts plus the midterm need 17.

\medskip
\noindent\textbf{What fitting would take.} About $1.3$ Parts per meeting, with
no review day and Midterm~2 on Nov~24. That is roughly two and a half times
the pace of matrix Part~1:

{\renewcommand{\arraystretch}{1.25}
\begin{center}
\small
\begin{tabular}{@{}l l @{\qquad} l l@{}}
\toprule
Oct 27 & Part 2 & Nov 19 & Part 11 (Gram--Schmidt read beforehand) \\
Oct 29 & Part 3 & Nov 24 & \textbf{Midterm 2} \\
Nov 3  & Parts 4 and 5 & Dec 1  & E1, and E2 begun \\
Nov 5  & Part 6 & Dec 3  & E2 and E3 \\
Nov 10 & Part 7, and Part 8 begun & Dec 8  & L1 and L2 \\
Nov 12 & Part 8 finished, and Part 9 & Dec 10 & L3 \\
Nov 17 & Part 10 &        & \\
\bottomrule
\end{tabular}
\end{center}}

\noindent\textbf{Levers}, since that is too fast:
\begin{enumerate}[label=(\alph*),leftmargin=2em]
  \item Assign lighter sections as reading: \S25 (rotations in $\R^{3}$),
        \S37 (principal minors), \S50 (Cramer's rule), \S54 beyond full
        column rank, E3's structure table, and L2's similarity invariants.
  \item Merge. Part~4's \S22 folds into Part~2, since each property is a
        one-line check. Part~7 can lean on Part~8 (determinant by pivots) and
        shrink to minors, cofactors, and Laplace expansion. E3's Gershgorin can
        close E2.
  \item Move L1~\S3 (the product rule from composition) forward beside \S13
        (D11). It costs little there and shortens the end of term.
  \item Move Midterm~2 earlier with a narrower scope, for example through
        Part~8, and examine least squares, the inverse, and $QR$ on the final.
        This changes ``Gram--Schmidt and $QR$ is the last topic for
        Midterm~2'', so it is the instructor's call.
  \item Trim the final's scope, for example by stating the SVD and its
        consequences without the full construction.
\end{enumerate}

\noindent\textbf{Homework.} Two of the five assignments remain after
Homework~3. A natural split is Homework~4 on Parts~2--6 (catalog Problems
4--12) and Homework~5 on Parts~7--11 (catalog 13--14). The eigen and
linear-maps units would then have no homework of their own.

\end{document}
